\section{Introduction}
\label{sec:intro}

An inherited state can affect both the descendants a cell produces and their
susceptibility to killing. Eliminating a marker along one lineage therefore
need not eliminate the family. Conversely, a declining expected population can
give a useful extinction guarantee, but a negative mean growth rate alone does
not quantify finite-founder survival. The appropriate control question concerns
the probability of leaving any surviving family under an explicitly priced
intervention.

Reaction-based models connect histone and DNA modifications to cellular memory
\citep{dodd2007,bruno,bruno2024}, and inherited non-genetic states are a
standard explanation for drug-tolerant persistence
\citep{balaban2004,sharma2010,shaffer2017,oren2021}. Inheritance and treatment
resistance have been studied in population models \citep{cassidy}; drug-induced
plasticity changes optimal dosing, and the chosen dose--response law can change
which schedules are favored \citep{gunnarsson}. Controlled branching processes
have a substantial stochastic-control theory \citep{claisse,kimmel2015}, and
tumour elimination probabilities have been used directly in treatment
objectives \citep{burroughs}. We do not claim priority for branching generating
functions, stochastic verification, or extinction-based treatment objectives.
Our contribution is a pair of explicit, checkable certificates --- one indexed
by the remaining budget, one uniform over the admitted action range --- their
founder-count and memory-size consequences, and a source-level realization that
retains correlated molecular inheritance.

\subsection*{Two obstructions, and a third}

Write $B$ for the total intervention exposure, meaning additional molecular
reaction rate integrated over time, and $v_{\max}$ for the largest instantaneous
rate the actuator can deliver. Three separate quantities can prevent
eradication.

\begin{itemize}
\item \textbf{Exposure.} Theorem~\ref{thm:main} shows that if a baseline
extinction supersolution $q$ obeys a componentwise inequality $Rq\le c(1-q)$
for the intervention generator, then survival from configuration $z$ is at
least $1-\prod_i[1-e^{-cB}(1-q_i)]^{z_i}$ for \emph{every} predictable policy
respecting the pathwise budget $B$. The exposure needed to reach risk $\delta$
from $n$ founders therefore grows like $c^{-1}\log(n/\delta)$.
\item \textbf{Amplitude.} Theorem~\ref{thm:amp} shows that if a single $q$ is a
supersolution for every admitted action --- equivalently, by affineness, at
both ends of the action range --- then survival is at least $1-\prod_iq_i^{z_i}$
for every policy at \emph{every} budget and horizon. This obstruction is not
removed by treating longer or paying more.
\item \textbf{Time.} Proposition~\ref{prop:time} shows that bounded death rates
impose a minimum duration for finite-time extinction, independently of the two
bounds above.
\end{itemize}

The three are logically independent, and they fail in different ways. Only the
first is a budget statement; only the second survives an unbounded course; only
the third constrains a fixed deadline.

\subsection*{What is new relative to the exposure floor alone}

The central technical device for the exposure floor is that no optimization over
a discretized population space is required. Given a baseline supersolution $q$,
one considers the entire segment $1-a(1-q)$, which is preserved by the
inequality a population barrier needs (Lemma~\ref{lem:segment}); setting $a$ by
the remaining budget gives a bound that approaches one with founder count at
every fixed finite exposure. A constant admissible action with a negative
weighted drift gives the matching logarithmic order
(Corollary~\ref{cor:order}), and an outward-rounded Taylor enclosure certifies a
sharper finite-horizon policy in the worked source
(Section~\ref{sec:pulse}).

The amplitude floor is what makes the memory-size question answerable. In the
molecular source of Section~\ref{sec:source} the erasure actuator has a
critical amplitude $v_c(N)$: below it the certificates exhibited here show that
\emph{no} policy eradicates, and above it a constant action does.
Section~\ref{sec:memory} brackets the corresponding
critical intrinsic erasure to six decimals for $2\le N\le8$ sites; it rises from
$.092750$ at two sites to $.450572$ at eight, a factor of $4.86$. In particular
the amplitude $.29$ used in the two-site worked example is certified
\emph{insufficient at any budget or duration} for $5\le N\le8$, with explicit
rational supersolutions bounding survival from a fully marked founder below by
$.170$, $.460$, $.613$ and $.701$ at $N=5,6,7,8$. This settles, negatively, a
question left open by a coarser mean-growth diagnostic, which had only shown
that a particular constant action ceases to contract.

Section~\ref{sec:memory} also shows that the deterioration is a property of this
actuator and not of inherited memory as such. A second actuator that adds death
on protected states has threshold at most $b-d_{\mathrm{prot}}$ for every $N$,
because beyond that value the expected total count decays at rate
$\min(\kappa-(b-d_{\mathrm{prot}}),\,d_{\mathrm{unprot}}-b)$ with unit weight
vector; the resulting sufficient exposure $1.45\log(n/\delta)$ is dimension-free
and lies within a factor $1.54$ of the exposure floor at $n=1$, $\delta=.01$.
Where the eraser alone needs a fourfold amplitude increase between two and eight
sites, a fixed added death of $.025$ restores mean contraction at six sites at
the erasure amplitude that by itself provably fails.

\subsection*{Scope}

The case study uses synthetic reaction and demographic rates. It is suitable for
exposing mathematical dependencies and for designing a low-density microcolony
test, not for prescribing a drug regimen. The exposure variable measures
additional molecular reaction rate integrated over time; it is not a drug mass
or a concentration, and Section~\ref{sec:pkpd} states exactly what is needed to
convert it. Throughout, ``eradication'' means the specified extinction event in
a process with no immigration, and ``survival'' means non-extinction of the
mathematical population, which is not a detection threshold, a clinical
recurrence or a progression endpoint.
